\chapter{DHCP Configuration}
\label{chap:dhcp}

\section{Background}
\label{sec:dhcp-background}
Every machine, from computers to servers, requires an IP address to be able to send and receive packets of data \cite{krause2019}. These IP addresses act like home addresses on a street, and whilst DNS is responsible for resolving those addresses into names, there are two ways of assigning IP addresses to machines: static addressing and DHCP \cite{krause2019}.

Static addressing is the process of configuring IP addresses on a system manually. It is a hands-on approach where the address is entered directly into the NIC settings of a particular device \cite{krause2019}. Despite its simplicity, this method does not scale in environments that host many devices, such as a small or large company with laptops, desktops, printers, servers, and other machines. Each of these devices needs an IP address, but assigning them manually becomes tedious at scale. The solution to this problem is the Dynamic Host Configuration Protocol (DHCP) \cite{krause2019}.

DHCP is a protocol that provides IP addressing information to a large number of devices automatically: a machine that is plugged into the network can obtain its address information from the DHCP server \cite{krause2019}. When a laptop or phone connects to a Wi-Fi router, a DHCP server assigns it the addressing information it needs to route traffic on the network. On public Wi-Fi the DHCP server runs on the router itself, whereas in businesses where Windows Server acts as the datacentre, DHCP services are most often hosted on one or more servers across the network \cite{krause2019}.

\section{Configuration Steps}
\label{sec:dhcp-steps}
The configuration is carried out in two stages: the post-install configuration of the DHCP role (\autoref{sec:dhcp-post}) and the creation of a DHCP scope (\autoref{sec:dhcp-scope}). \autoref{tab:dhcp-config} lists the values used throughout this chapter. If the DHCP console cannot be used, or something goes wrong, see \autoref{sec:dhcp-issues}.

\begin{table}[H]
  \centering
  \small\sffamily
  \renewcommand{\arraystretch}{1.3}
  \begin{tabular}{@{}p{5.6cm}p{8.6cm}@{}}
    \toprule
    \textbf{Setting} & \textbf{Value used} \\
    \midrule
    Scope name & GROUP 1-SCOPE \\
    Address pool & 192.168.0.100 -- 192.168.0.254 \\
    Subnet mask & 255.255.255.0 (/24) \\
    Exclusions and delay & None \\
    Lease duration & 8 days (default) \\
    Default gateway (router) & 192.168.0.1 \\
    Parent domain & OPS.local \\
    DNS servers & 192.168.0.5 and 8.8.8.8 \\
    WINS servers & None \\
    Scope state & Activated \\
    Authorisation credentials & \texttt{OPS\textbackslash Administrator} \\
    \bottomrule
  \end{tabular}
  \caption{Key configuration values used in this chapter}\label{tab:dhcp-config}
\end{table}

%--------------------------------------------------------------
\subsection{DHCP Post-Install Configuration}
\label{sec:dhcp-post}
Like other server roles, DHCP must first be added to the server during the installation phase. After the DHCP role has been added, it is not yet properly configured.

The DHCP server can give clients two kinds of access groups: DHCP Administrators and DHCP Users \cite{dhcp-groups}. DHCP Users have read-only access to the server, which allows them to view, but not modify, server data, including the DHCP server configuration, registry keys, DHCP log files, and the database. They cannot create scopes, modify option values, or create reservations or exclusion ranges. DHCP Administrators are able to view and modify data at the DHCP server. They can create and delete scopes, add reservations, change option values, create superscopes, or perform any other activity needed to administer the DHCP server. Despite this control, they do not have unlimited administrative rights, because they only have rights on the local computer \cite{dhcp-groups}.

These two groups are local to the server, which means they cannot authorise DHCP servers; this can only be done by Domain Admins \cite{dhcp-groups}. Authorising the server is needed because DHCP communication typically occurs before any user or computer authentication, and because DHCP is based on IP, an unknown DHCP server could hand invalid information to clients. Authorisation prevents this \cite{msnoob-dhcp}.

The authorisation step in the wizard requires the credentials of an account that is allowed to authorise the server in Active Directory (\autoref{fig:dhcp-01} and \autoref{fig:dhcp-02}). So far the only credential set was \texttt{OPS\textbackslash Administrator}, and any user who attempts to connect to the server in future must use this username. The wizard confirms the result in \autoref{fig:dhcp-03}.

\shot{dhcp_01.png}{Description of the DHCP Post-Install Configuration Wizard}{fig:dhcp-01}
\shot{dhcp_02.png}{Authorisation of the DHCP Post-Install Configuration Wizard}{fig:dhcp-02}
\shot{dhcp_03.png}{Post-Install Configuration Completed}{fig:dhcp-03}

\begin{callout}{Restart the DHCP service}
The wizard asks for the DHCP server service to be restarted so that the new security groups become effective. Do this before creating the scope.
\end{callout}

%--------------------------------------------------------------
\subsection{The DHCP Scope}
\label{sec:dhcp-scope}
For a DHCP server to hand out IP addresses, it needs to know which IP addresses, subnet mask, default gateway and DNS server addresses are valid within the network, so that it can package that information and hand it out to the computers that request it. This package of information is called a DHCP scope \cite{krause2019}.

To begin defining the scope, the DHCP management tool was launched from the Tools menu in Server Manager (\autoref{fig:dhcp-04}). This opens a window for DHCP with the name of the device that the server is running on (\autoref{fig:dhcp-05}).

\shot{dhcp_04.png}{DHCP Selected on the Tools Dropdown Menu}{fig:dhcp-04}
\shot{dhcp_05.png}{DHCP Window}{fig:dhcp-05}

After opening the DHCP window, the server was expanded to show two protocols: IPv4 and IPv6. IPv4 was selected, which displayed information about adding a scope, along with the server options, policies and filters. IPv4 did not have a scope yet, so one needed to be configured. IPv4 was right-clicked to open the menu containing the option to add a new scope (\autoref{fig:dhcp-06}). Choosing this option opened the New Scope Wizard (\autoref{fig:dhcp-07}). After clicking Next, the basics of the scope were decided. The name given was GROUP 1-SCOPE (\autoref{fig:dhcp-08}).

\shot{dhcp_06.png}{Adding a Scope}{fig:dhcp-06}
\shot{dhcp_07.png}{New Scope Wizard}{fig:dhcp-07}
\shot{dhcp_08.png}{Group One Scope Made}{fig:dhcp-08}

There exist 4.3 billion possible IPv4 addresses, of which roughly 3.7 billion are available for use \cite{ipme-ipv4}. A scope does not need to accommodate billions of addresses; it just needs a few hundred or a few dozen. When configuring it, the range of IP addresses had to be decided. The addresses that the scope will hand out are 192.168.0.100 to 192.168.0.254 (\autoref{fig:dhcp-10}), and this is the address pool. No exclusions or delays were made in the scope (\autoref{fig:dhcp-09}).

Alongside the handed-out IP addresses, the scope also has the option to configure the amount of time that a client keeps an assigned IP address before renewal, ensuring efficient management of IP addresses \cite{rapidseedbox-lease}. This period of time is called the lease duration. It was left at the default of 8 days (\autoref{fig:dhcp-11}).

\shot{dhcp_09.png}{No Exclusions Made}{fig:dhcp-09}
\shot{dhcp_10.png}{Setting the IP Address Range}{fig:dhcp-10}
\shot{dhcp_11.png}{Lease Duration of 8 Days Set}{fig:dhcp-11}

When clients use the server, access to the internet is an important feature. In order for the clients to be able to access the internet, they need to be able to talk to devices and networks outside of the DHCP subnet. A default gateway allows for this; without it, a client can only reach devices and networks inside its own subnet. The default gateway was given the IP address 192.168.0.1 (\autoref{fig:dhcp-12} and \autoref{fig:dhcp-13}).

\shot{dhcp_12.png}{Default Gateway Configured (1)}{fig:dhcp-12}
\shot{dhcp_13.png}{Default Gateway Configured (2)}{fig:dhcp-13}

After the default gateway, the next step was setting up the DNS server. The DHCP server hands out IP addresses and more, so that a client can use the network. An IP address allows a device to communicate with other hosts, but IP addresses are inconvenient to memorise or type. The DNS server and domain name exist so that devices can reach hosts by names that resolve to IP addresses \cite{msft-dhcp-dns}. The domain name was set to OPS.local, and two DNS server addresses, 192.168.0.5 and 8.8.8.8, were configured (\autoref{fig:dhcp-14}). This means that, in order to access the server, a client simply needs to type in the domain name, which makes it more accessible for non-domain-joined machines.

\shot{dhcp_14.png}{DNS Server and Domain Name}{fig:dhcp-14}

After the DNS server and domain name were configured, there was an option for WINS server configuration. WINS is a part of the scope that resolves NetBIOS computer names to IP addresses. It is needed to support pre-Windows~2000 systems and older applications that use NetBIOS over TCP/IP \cite{msft-wins}. None of the systems planned to be connected to the network qualify as pre-Windows~2000 systems, therefore WINS is not needed and was left vacant (\autoref{fig:dhcp-15}).

\shot{dhcp_15.png}{WINS Server Configuration Left Vacant}{fig:dhcp-15}

All necessary parts of the scope have now been configured. The only step left was the activation of the scope (\autoref{fig:dhcp-16}). Upon activating the scope, the wizard presents a completion display, and the Finish button was clicked to close it (\autoref{fig:dhcp-17}). The scope is now activated and ready for testing. It will hand out IP addresses from 192.168.0.100 to 192.168.0.254. At the time of creation there were no address leases yet (\autoref{fig:dhcp-18}).

\shot{dhcp_16.png}{Activation of the Scope}{fig:dhcp-16}
\shot{dhcp_17.png}{Finishing the New Scope Wizard}{fig:dhcp-17}
\shot{dhcp_18.png}{New DHCP Scope Created}{fig:dhcp-18}

%--------------------------------------------------------------
\section{Issues Encountered for DHCP Server Setup and Alternative Paths}
\label{sec:dhcp-issues}

\issue{DHCP Console Stops Functioning}
{The DHCP console can stop working, for example after a domain name change or after a new Ethernet adapter is added for internet access, so a new scope cannot be configured through the window.}
{Windows PowerShell is still available on Windows Server 2019 and bypasses surface-level console problems. The same configuration can be done with PowerShell commands (see \autoref{sec:dhcp-ps} and \autoref{sec:dhcp-psscope}). They can also be used when you simply prefer PowerShell.}

%--------------------------------------------------------------
\subsection{Use of PowerShell in the Event of DHCP Console Failure}
\label{sec:dhcp-ps}
A failing DHCP console does not mean that new DHCP scopes cannot be configured or removed. PowerShell commands allow DHCP scopes to be configured and modified with a similar degree of control as the console. Some of the important commands for configuring DHCP IPv4 scopes are listed in \autoref{tab:dhcp-cmdlets} \cite{msft-dhcpserver-module}, and examples of their output are shown in \autoref{fig:dhcp-19} to \autoref{fig:dhcp-22}.

\begin{table}[H]
  \centering
  \small\sffamily
  \renewcommand{\arraystretch}{1.25}
  \begin{tabular}{@{}p{7.2cm}p{7cm}@{}}
    \toprule
    \textbf{Command} & \textbf{Purpose} \\
    \midrule
    \texttt{\footnotesize Add-DhcpServerv4Scope} & Add scopes \\
    \texttt{\footnotesize Get-DhcpServerv4Scope} & View scopes \\
    \texttt{\footnotesize Set-DhcpServerv4Scope} & Modify scopes \\
    \texttt{\footnotesize Remove-DhcpServerv4Scope} & Delete scopes \\
    \texttt{\footnotesize Get-DhcpServerv4ScopeStatistics} & Show usage statistics \\
    \texttt{\footnotesize Add-DhcpServerv4Lease} & Add a lease \\
    \texttt{\footnotesize Get-DhcpServerv4Lease} & View leases \\
    \texttt{\footnotesize Remove-DhcpServerv4Lease} & Remove a lease \\
    \texttt{\footnotesize Add-DhcpServerv4Reservation} & Add reservations \\
    \texttt{\footnotesize Get-DhcpServerv4Reservation} & View reservations \\
    \texttt{\footnotesize Set-DhcpServerv4Reservation} & Modify reservations \\
    \texttt{\footnotesize Remove-DhcpServerv4Reservation} & Delete a reservation \\
    \bottomrule
  \end{tabular}
  \caption{PowerShell commands for managing DHCP IPv4 scopes}\label{tab:dhcp-cmdlets}
\end{table}

\shot{dhcp_19.png}{Example of Accessing DHCP Through PowerShell}{fig:dhcp-19}
\shot{dhcp_20.png}{Example of Adding a DHCP Reservation}{fig:dhcp-20}
\shot{dhcp_21.png}{Example of Verification of a DHCP Reservation}{fig:dhcp-21}
\shot{dhcp_22.png}{Example of IP Address Tracking}{fig:dhcp-22}

%--------------------------------------------------------------
\subsection{Using Windows PowerShell to Create a DHCP Scope}
\label{sec:dhcp-psscope}
To create a DHCP scope using Windows PowerShell, it must be opened as Administrator \cite{sherry-dhcp-scope}. Considering that scopes have most likely been made already, a command must first be run to check for existing scopes.

The command \texttt{Get-DhcpServerv4Scope} lists every scope on the DHCP server that was created using the DHCP console (\autoref{fig:dhcp-23}) \cite{sherry-dhcp-scope}. If there are no preexisting scopes, nothing is shown.

\shot{dhcp_23.png}{Example of the Get-DhcpServerv4Scope Command in Use}{fig:dhcp-23}

Creating a scope in PowerShell would normally involve a lot of typing, as the \texttt{Add-DhcpServerv4Scope} command has many parameters. A conventional use of the full command would look like the following:

\begin{lstlisting}
Add-DhcpServerv4Scope `
    -Name "Test Scope" `
    -Description "Scope for DHCP testing" `
    -StartRange 192.168.10.100 `
    -EndRange 192.168.10.200 `
    -SubnetMask 255.255.255.0 `
    -LeaseDuration 0.08:00:00 `
    -State Active `
    -PassThru
\end{lstlisting}

This can be tedious, so splatting can be used to make passing parameters easier \cite{listek-splatting}. PowerShell splatting is when a hashtable is used to create a parameter set that is passed to the cmdlet, which then creates the DHCP scope according to the given parameters. For example, if the following were entered in the PowerShell console:

\begin{lstlisting}
$HashArgs = @{
    'Name'          = 'Head Office Wireless Scope'
    # This sets the name of the wireless scope
    'Description'   = 'Wireless Clients'
    # This is the human-readable description of the scope
    'StartRange'    = '10.1.0.100'
    # Specifies the starting IP address in the scope
    'EndRange'      = '10.1.0.200'
    # Specifies the end IP address in the scope
    'SubnetMask'    = '255.255.255.0'
    # Specifies the subnet mask of the scope
    'State'         = 'Active'
    # Activates the scope
    'LeaseDuration' = '1.00:00:00'
    # Specifies the length of the lease duration
}
Add-DhcpServerv4Scope @HashArgs
\end{lstlisting}

the hashtable \texttt{\$HashArgs} is given the DHCP scope parameters. When the \texttt{Add-\allowbreak DhcpServerv4Scope} command is entered, a new scope named Head Office Wireless Scope is made, with an address pool between 10.1.0.100 and 10.1.0.200 and a subnet mask of 255.255.255.0 \cite{sherry-dhcp-scope}. The lease duration is one day and the scope is made active. To check that the scope was made, run \texttt{Get-DhcpServerv4Scope} again. A way to make the scopes easier to read is to add \texttt{| Format-List} at the end of the command \cite{sherry-dhcp-scope}.

%--------------------------------------------------------------
\section{Testing}
\label{sec:dhcp-testing}

\subsection{Checking the Server}
\label{sec:dhcp-checkserver}
To test the operation of a DHCP server, there must be confirmation that data packets are travelling between the server and client devices. To start, the server is checked in a PowerShell window that is run as Administrator, using the command:

\begin{lstlisting}
Get-Service -Name DHCPServer
\end{lstlisting}

The Status column must display \emph{Running}. If the command shows the service as \emph{Stopped}, the server cannot answer any client requests. The reason may be that the service has not started or that the server has not been authorised. If it needs to be started, run

\begin{lstlisting}
Start-Service -Name DHCPServer
\end{lstlisting}

and then run the \texttt{Get-Service} command again to confirm that it is running. Another probable reason is that the server needs authorisation. The command \texttt{Get-DhcpServerInDC} verifies whether the DHCP server is authorised in Active Directory \cite{msft-dhcp-quickstart}. For a server that is not yet listed, it must first be added to the directory with:

\begin{lstlisting}
Add-DhcpServerInDC -DnsName <server FQDN> -IPAddress <server IP address>
\end{lstlisting}

A few other checks are to confirm the configuration and the scope usage, seeing how many addresses are in use and how many are free \cite{msft-dhcp-quickstart}:

\begin{lstlisting}
Get-DhcpServerv4ScopeStatistics
Get-DhcpServerv4Scope
\end{lstlisting}

\subsection{Testing from a Client and Verifying on the Server}
\label{sec:dhcp-testclient}
To test the server from a client, the client has to request a lease. On a Windows client, open Command Prompt and run \texttt{ipconfig /release}, then \texttt{ipconfig /renew}, then \texttt{ipconfig /all}. The first command releases the DHCP client lease, the second renews it, and \texttt{ipconfig /all} (or just \texttt{ipconfig}) verifies the DHCP client lease (\autoref{fig:dhcp-24}) \cite{msft-dhcp-renew}.

\shot{dhcp_24.png}{\texttt{ipconfig /all} Command in Command Prompt}{fig:dhcp-24}

The address, subnet mask, gateway, DNS servers and DHCP server line must all match the scope (\autoref{tab:dhcp-config}).

\issue{Client Has Incorrect or Missing Options}
{The client lease shows incorrect or missing options. This can happen because the network configuration is not working correctly, because the client cannot support the missing options, or because the server is not configured to hand them out \cite{msft-dhcp-clients}.}
{Change the list of values for the missing or incorrect options at the DHCP scope in use. You might have to configure the options as a Server Option at the DHCP server, remove it, and set the correct value in the Scope Options of the DHCP scope that is used by the client \cite{msft-dhcp-clients}.}

After the lease is recorded, run the following command to confirm that clients are receiving addresses. The cmdlet also has a means of identifying bad leases with \texttt{-BadLeases}, and \texttt{-AllLeases} specifies all IPv4 address leases \cite{msft-get-lease}:

\begin{lstlisting}
Get-DhcpServerv4Lease -ScopeId <scope ID>
\end{lstlisting}

The server logs and lease records are kept in two places: major DHCP service events are tracked in Event Viewer, but individual lease events go to separate DHCP log files. By default, these are in the directory \path{C:\Windows\System32\dhcp}, with one file per weekday \cite{garn-dhcp-logs}.
